% Transcription — fonds Alexandre Grothendieck, shelfmark 156-4, batch 2 (pages 21-40).
% Established from the facsimile archives/batches/156-4/batch-02.pdf (a copy of
% 156-4.pdf fetched through the project's relay; no cover sheet in the batch, so
% PDF page k = folder page 20+k), per the skill transcribe-grothendieck. The
% folder is eighty-eight pages in five batches; this is the second. Batch 1
% (pages 1-20) was transcribed in parallel and read once this batch was
% drafted: page 21 continues its page 20 (L_X, L̂_X, L_F, L̂_F, a cancelled
% « Contr_3 »), and its notation is kept — \mathcal{L}, \mathcal{M} for
% lieux and multistrates, \mathcal{F} for a figure on pages 22-23. Batch 1
% reads his axiom label « Contr »; on these pages the ink reads « Cont »
% (p. 21 clearly), transcribed as read and noted on page 21.
%
% Pass: Opus 5.5 (claude-opus-5-5), 2026-09-26 — first pass, unchecked against the pages by a human.
%
% Pages skipped: none. Pages summarised: none. All twenty leaves are his
% mathematics, in ink, fast drafting throughout.
%
% His pagination 21-40 runs at the top of the leaves and coincides with the
% archivists'. One date in his hand: « 12. Juin » (underlined) opening a
% paragraph on page 31, which begins « Depuis hier j'ai pas mal réfléchi… »;
% it is transcribed where it stands and noted. No other date on these pages.
% His numbered runs are the axioms « Cont » (contrées), numbered and
% renumbered as he goes: Cont 3 (p. 21), Cont 4 (p. 22), Cont 5 (p. 24); then,
% after « On reprend tout ! » on page 25, a fresh run Cont 1 (p. 25), Cont 2
% (p. 26), Cont 3 (p. 27), Cont 4 (p. 29), Cont'_4 (pp. 30-31); then, after the
% change of framework on pages 31-32, a third run: Cont 1 and the axioms
% « Fig 1 » to « Fig 4 » (pp. 33-34), Cont'_3 and Cont 3 (p. 35), Cont 4
% (p. 35), Cont 5 and Cont 6 (p. 39). Several of these numbers are written
% over others; each is noted where it occurs. Page 39's margin cross-refers
% to « p. 33 ».
%
% What the batch is about. It is the first draft (« première mouture ») of
% chapter IV, Analysis situs, and these pages try three times to axiomatise a
% « contrée ». Pages 21-24: in the framework of the preceding batch
% (multiplicities M with the refinement relation ≪, lieux L, the set L_X of
% lieux x ≪ X), he defines the open stratum L°_X = L_X ∖ ∪_{X_i<X} L_{X_i} and
% asks that it be non-empty (Cont 3) and that the L° of distinct compatible
% multiplicities be disjoint (Cont 4); a Scholie (p. 23) that for a figure F
% the L°_X partition L_F, X ↦ L_X is injective, and X ≤ Y ⟺ L_X ⊂ L_Y ⟺
% L°_X ⊂ L_Y ⟺ L°_X ∩ L_Y ≠ ∅; a converse Cont 5 (p. 24) recovering ≪ and
% compatibility from the sets L̂_X. Page 25, « On reprend tout ! »: a contrée
% is a set L of lieux with a set 𝔉 ⊂ P(P(L)) of « figures », whose members
% are the (closed) strata; support |F|; Cont 1 (empty and point figures),
% Cont 2 (open strata X° = X ∖ ∂X non-empty, X ∩ Y a union of strata),
% Cont 3 (closed parts of a figure are sub-figures, F' ≤ F), compatibility of
% two figures (four equivalent conditions, p. 28), « configurations », Cont 4
% (p. 29) and « contrées de points » (ensemblistes), disjoint lieux and
% C-disjoint parts (p. 30). Pages 31-32 (12 June): reflecting on topological
% pseudo-discs, he concludes that lieux must carry extra data and that a
% contrée should rather be C = (M, ≪, R_M, F); he resolves to fix a
% geometric terminology first and axiomatise later. Pages 33-40: the axioms
% in that framework — Cont 1 (X ≤ Y ⟺ X ≪ Y and X compatible with Y is an
% order), Fig 1-4 (elementary figures M_{≤X}, closed parts, union of
% compatible figures), « contrées finitistes » (p. 34), finiteness of
% combinatorial and geometric dimension (Cont'_3, Cont 3), lieux as minimal
% elements, the « ombre » L_X of X on L and the open ombre (p. 35), a Scholie
% that the ombres of a figure form an admissible family of parts of L (p. 36),
% sub-contrées (p. 37), disjoint multistrates, refinement of figures F ⋘ G and
% « sub-figure ⟺ refines and compatible » (p. 38), Cont 5 and Cont 6 on
% disjointness (p. 39), and « contrées induites », the residual contrée
% C_{∁F}, and the definition of a subdivision of a figure (p. 40).
%
% Reading hazards fixed once here. His script M (multistrates, the word batch
% 1 p. 19 adopts; « multiplicité » on p. 21) is set \mathcal{M}, his script L
% (lieux) \mathcal{L}, a single figure on pages 22-23 \mathcal{F} (as batch 1
% p. 20), the set of figures from page 25 on \mathfrak{F}, the power set
% \mathfrak{P}. Batch 1 sets multistrates as a bold X where he bars it; no
% bar was seen on these pages and plain X is kept. Refinement is ≪ (\ll); on page 38 the relation
% between figures is a triple ⋘ (\lll). His tilde X̃ is set \widetilde{X}. His
% underlining is set \emph. His circled (a), (b) on page 39 are set in
% parentheses with a note. Long slanted left-margin notes on pages 24, 26, 28,
% 30, 32, 33, 34, 35, 36, 38, 39, 40 were read only in part; they are given as
% \marginal where something could be read, otherwise noted.
%
% Readings flagged and not settled: most connective prose on pages 30 and
% 32-33; the heading of page 38 (« Figures » struck, an insertion above);
% the superscript on M in « M^♯_≪ » (p. 40); « Cont » numbers written over
% other numbers (pp. 30, 34, 35). Variances on the page and not repaired:
% « des Y ∈ X » for Y ≤ X (p. 33); « F ∪ C » for F ∪ G (p. 28).

\documentclass[11pt,a4paper]{article}
\input{../preamble/grothendieck}

\folder{156-4}
\batch{2}
\pages{21}{40}
\dating{1986}
\watermark{Édition de démonstration}
\foldertitle{[Chapitre] IV. Analysis situs (première mouture) : notes manuscrites (10/06/1986)}

\begin{document}

\page{21}
\note{la page continue la fin de la page 20 (batch 1), où sont posés $\mathcal{L}_X$, $\widehat{\mathcal{L}}_X$, $\mathcal{L}_{\mathcal{F}}$, $\widehat{\mathcal{L}}_{\mathcal{F}}$ et où un « Contr$_3$ » est commencé puis barré ; l'abréviation lue ici « Cont » est vraisemblablement la même}
\struck{Soit $X$ une multiplicité.} L'application $X_i \mapsto \mathcal{L}_{X_i}$, $\mathcal{M}_{\leq X} \to \mathfrak{P}(\mathcal{L})$ est bien sûr croissante ($X_i \leq X_j \Rightarrow \mathcal{L}_{X_i} \subset \mathcal{L}_{X_j}$). On aimerait qu'elle le soit \emph{strictement}, i.e. que $X_i < X_j \Rightarrow \mathcal{L}_{X_i} \subsetneq \mathcal{L}_{X_j}$, i.e. $\exists$ lieu $x$ tel que $x \ll X_j$ et $x \not\ll X_i$.
\note{les indices $i$, $j$ sont plusieurs fois récrits l'un sur l'autre}

Soit $\mathcal{L}^{\circ}_X = \mathcal{L}_X \smallsetminus \bigcup_{X_i < X} \mathcal{L}_{X_i}$. Condition plus forte :

\emph{Cont 3} \quad $\forall X \in \mathcal{M}$, $\mathcal{L}^{\circ}_X \neq \emptyset$ — i.e. $\forall X \in \mathcal{M}$, $\exists x \in \mathcal{L}$, tel que $x \ll X$, et $x \not\ll X_i$ pour tout $X_i < X$.
\note{« Cont 3 » est écrit dans la marge gauche, souligné, en regard de la condition}

\struck{Ainsi, si $X \in \mathcal{M}$, l'application $X_i \mapsto \mathcal{L}^{\circ}_{X_i}$, $\mathcal{M}_{\leq X} \to \mathfrak{P}(\mathcal{L})$ \ill{}. Soient $X$ et $Y$ \ill{} multiplicités. Supposons \ill{} compatibles, \ill{}}
\note{ces lignes sont barrées de traits horizontaux et de deux longs traits obliques ; seul le début se lit}

Soit $X \in \mathcal{M}$, considérons les $\mathcal{L}^{\circ}_{X_i}$, $X_i \leq X$. Je dis que leur réunion (qui est $\subset \mathcal{L}_X$) est égale à $\mathcal{L}_X$. En effet, soit $x \in \mathcal{L}_X$, donc \struck{$\exists$} \ill{}, $x \in \mathcal{L}_{X_i}$, \struck{\ill{}} et \struck{\ill{}} il y a un $X_i$ minimal parmi les

\page{22}
$X_i$ tels que $x \in \mathcal{L}_{X_i}$ (\ill{}) donc $x \in \mathcal{L}^{\circ}_{X_i}$. De façon précise, \ill{} : $x \in \mathcal{L}^{\circ}_{X_i} \Leftrightarrow X_i$ est minimal parmi les $Y \in \mathcal{M}_{\leq X}$ tels que $x \in \mathcal{L}_Y$.

J'aimerais que cet $X_i$ soit \emph{unique}, i.e. qu'il soit \uncertain{même} plus petit élément parmi les $Y$ envisagés, donc $X_i, X_j \leq X$, $X_i \neq X_j \Rightarrow \mathcal{L}^{\circ}_{X_i} \cap \mathcal{L}^{\circ}_{X_j} = \emptyset$. Donc que $\mathcal{L}_X$ soit réunion \emph{disjointe} des $\mathcal{L}^{\circ}_{X_i}$ pour $X_i \leq X$. Plus généralement, on aimerait que l'analogue \uncertain{soit vrai} pour \struck{tt} une figure $\mathcal{F} \subset \mathcal{M}$, \struck{donc} que $\mathcal{L}_{\mathcal{F}}$ soit réunion \emph{disjointe} des \struck{\ill{}} $\mathcal{L}^{\circ}_{X_i}$ pour $X_i \in \mathcal{F}$ (alors qu'on sait \add{déjà} que c'est la réunion des $\mathcal{L}^{\circ}_{X_i}$) : on a encore que $\forall x \in \mathcal{L}_{\mathcal{F}}$, $\exists$ plus \emph{petit} $X \in \mathcal{F}$, tel que $x \in \mathcal{L}_X$ \struck{\ill{}}.

\emph{Cont 4} \quad Si $X$, $Y$ sont deux multiplicités \emph{compatibles} (\struck{définissant une} \add{\ill{} dans une} figure) et $X \neq Y$, alors $\mathcal{L}^{\circ}_X \cap \mathcal{L}^{\circ}_Y = \emptyset$. (En d'autres termes, si $x \in \mathcal{L}_X \cap \mathcal{L}_Y$, alors $\exists$ \struck{$Z'$} $X' \leq X$ tel que $x \in \mathcal{L}_{X'}$ et $Y' \leq Y$ tel que $x \in \mathcal{L}_{Y'}$.) \struck{D}

\marginal{(ceci équiv. à $\mathcal{L}_X \cap \mathcal{L}_Y = \bigcup_{Z \leq X,\, Z \leq Y} \mathcal{L}_Z$ \ill{})}

\struck{\ill{} Mais on peut \ill{} que : si $X, Y \in \mathcal{M}$, $\forall X' \leq X$, $Y' \leq Y$, avec $X' \neq Y'$, on ait $\mathcal{L}^{\circ}_{X'} \cap \mathcal{L}^{\circ}_{Y'} = \emptyset$, alors $X'$ et $Y'$ sont compatibles}
\note{passage hachuré, lu en partie seulement}

\page{23}
En utilisant \emph{Cont 3}, \emph{Cont 4}, on trouve le

\emph{Scholie.} Soit $\mathcal{F}$ ($\subset \mathcal{M}$) une figure. \struck{L'application}
\begin{enumerate}
  \item[a)] Les \struck{\ill{}} \add{parties de $\mathcal{L}$} (« strates ouvertes de $\mathcal{F}$ ») $\mathcal{L}^{\circ}_X$ ($X \in \mathcal{F}$) sont non vides et mutuellement disjointes, elles forment une partition de $\mathcal{L}_{\mathcal{F}} = \bigcup_{X \in \mathcal{F}} \mathcal{L}_X = \bigcup_{X \in \mathcal{F}} \mathcal{L}^{\circ}_X$. \struck{b) L'application} À fortiori, l'application $X \mapsto \mathcal{L}^{\circ}_X$, $\mathcal{F} \to \mathfrak{P}(\mathcal{L})$ [est injective].
  \item[b)] L'application $X \mapsto \mathcal{L}_X$, $\mathcal{F} \to \mathfrak{P}(\mathcal{L})$ est injective. On a, pour $X, Y \in \mathcal{F}$, équivalence des conditions suivantes : (i) $X \leq Y$ \quad (ii) $\mathcal{L}_X \subset \mathcal{L}_Y$ \quad (iii) $\mathcal{L}^{\circ}_X \subset \mathcal{L}_Y$ \quad (iv) $\mathcal{L}^{\circ}_X \cap \mathcal{L}_Y \neq \emptyset$.
  \item[c)] Si $X, Y \in \mathcal{F}$, on a $\mathcal{L}_X \cap \mathcal{L}_Y = \bigcup_{Z \in \mathcal{F},\, Z \leq X,\, Z \leq Y} \mathcal{L}_Z$.
\end{enumerate}
\marginal{dans ces équivalences, \ill{} dit que $X$, $Y$ sont des multiplicités \emph{compatibles}. — Idem pour c)}
\note{la remarque de droite est encadrée, en regard de b) ; « Idem pour c) » est en regard de c). Une note oblique en marge gauche, à hauteur de a), n'est lue qu'en fragments (« Mais bien sûr… », « multiplicités », « injective »)}

\emph{Dém.} de b). \ill{} (i) $\Rightarrow$ (ii) $\Rightarrow$ (iii) $\Rightarrow$ (iv) triviales, prouvons (iv) $\Rightarrow$ (i). Comme $\mathcal{L}_Y = \bigcup_{Y_i \leq Y} \mathcal{L}^{\circ}_{Y_i}$, il existe $Y_i$ avec $\mathcal{L}^{\circ}_X \cap \mathcal{L}^{\circ}_{Y_i} \neq \emptyset$, par a) on a $X = Y_i$ donc $X \leq Y$, cqfd.
\note{l'indice $Y_i$ est écrit sur un $Z_i$ barré}

\page{24}
\marginal{\ill{}}
\note{note oblique dans le coin supérieur gauche, illisible}
Mais on en arrive aux réciproques.

\emph{Cont 5} \quad Soient $X, Y \in \mathcal{M}$ deux multiplicités, et considérons \struck{\ill{}} $\widehat{\mathcal{L}}_X$, $\widehat{\mathcal{L}}_Y \subset \mathfrak{P}(\mathcal{L})$ comme ensembles de parties de $\mathcal{L}$. Ils permettent, \uncertain{on le voit}, de récupérer les relations \ill{} correspondantes :
\begin{enumerate}
  \item[a)] Pour que $X \ll Y$, il faut et il suffit \struck{que $\forall A \in \widehat{\mathcal{L}}_X$} que toute strate ouverte de $X$ soit contenue dans une strate ouverte de $Y$.
  \item[b)] Pour que $X$ et $Y$ soient compatibles, il faut et il suffit que pour toute strate ouverte de $X$ et toute strate ouverte de $Y$, \struck{on ait} celles-ci soient disjointes ou identiques.
\end{enumerate}

\emph{Cor.} a) $\widehat{\mathcal{L}}_X = \widehat{\mathcal{L}}_Y \Leftrightarrow X = Y$ \quad b) $\widehat{\mathcal{L}}_X \subset \widehat{\mathcal{L}}_Y$ \struck{\ill{}} fermé dans $\widehat{\mathcal{L}}_Y \Leftrightarrow X \leq Y$.
\note{les signes $=$ de a) sont surchargés et douteux}

\emph{Scholie.} On récupère $\mathcal{M}$, $\ll$, $R$ (compatibilité) \ill{} en termes de $\mathcal{L}$ et de l'ens. $\widehat{\mathcal{M}}$ des $\widehat{\mathcal{L}}_X$ ;
\[
  \widehat{\mathcal{M}} = \operatorname{Im}\bigl(\underbrace{\mathcal{M} \xrightarrow{X \mapsto \widehat{\mathcal{L}}_X} \mathfrak{P}(\mathfrak{P}(\mathcal{L}))}_{\text{application injective}}\bigr) \in \mathfrak{P}(\mathfrak{P}(\mathfrak{P}(\mathcal{L})))
\]
\note{une note oblique dans la marge gauche, au milieu de la page, n'est lue qu'en fragments}

\section{On reprend tout !}
\note{ces mots, soulignés, sont en haut à droite de la page 25, avec une grande lettre $C$ au-dessus}

\page{25}
\emph{Définition-scholie.} Une \emph{contrée} $C$ est définie par un ensemble $L$ (l'ens. des \emph{lieux} de la contrée), et la donnée de certains ensembles de parties de $L$ (de la contrée), appelées les \emph{figures} de $C$. L'ens. $\mathfrak{F}$ des figures de $C$ est donc une partie de $\mathfrak{P}(\mathfrak{P}(L))$ (donc $\mathfrak{F} \in \mathfrak{P}(\mathfrak{P}(\mathfrak{P}(L)))$). Si $F$ est une figure, \struck{ensemble} les éléments \add{(qui sont des parties de $L$)} s'appellent les \emph{strates} (\emph{fermées}) de la figure $F$. La réunion des strates \struck{fermées} de $F$ est appelée le \emph{support} de $F$, et notée $|F|$. On pose des \emph{axiomes des contrées}, qui vont être dégagés par la suite.

\emph{Cont 1} \struck{Axiome} a) La partie vide de $\mathfrak{P}(L)$ est une figure, appelée « \emph{figure vide} ». C'est donc la figure unique qui n'a pas de strates.

b) $\forall x \in L$, \struck{la} \add{il y a une} figure \ill{} \struck{les} strates unique $\{x\}$. On l'appelle figure \emph{ponctuelle} définie par $x$.
\note{un double trait vertical dans la marge gauche en regard de b) ; une note oblique dans la marge, illisible}

\page{26}
\emph{Cont 2} \quad Soit $F$ une figure. \struck{Alors} Pour tout $X \in F$, \struck{l'ensemble} soit
\[
  \partial X = \bigcup_{Y \in F,\ Y \subsetneq X} Y, \qquad X^{\circ} = X \smallsetminus \partial X
\]
(strate \emph{ouverte} de $F$ définie par $X$). Alors on a $X^{\circ} \neq \emptyset$. De plus, si $X, Y \in F$,
\[
  X \cap Y = \bigcup_{Z \in F,\ Z \leq X, Y} Z
\]
i.e. \struck{tt} pour tt \struck{pt} lieu qui appartient aux deux strates fermées, \struck{on a $\exists Z \in F$} il existe une strate f. $Z \subset X \cap Y$ \struck{telle que} qui contient $x$.
\note{« $X^{\circ} \neq \emptyset$ » est encadré}
\marginal{Les \ill{} \ill{} ont été dites « admissibles »}
\note{lecture de la note oblique de la marge gauche en partie seulement}

\emph{Corollaire} \struck{Simplification} \add{\uncertain{Justification}} : $X \mapsto X^{\circ}$ induit une bijection de $F$ avec l'ens. des classes d'une relation d'équivalence dans $|F|$, i.e. pour $X, Y \in F$, \struck{tel que} \struck{$X \neq Y$, on a $X^{\circ} \cap Y^{\circ} = \emptyset$} $X \neq Y \Leftrightarrow X^{\circ} \cap Y^{\circ} = \emptyset$ (ce qui \ill{} fait : \emph{Cont 2}).

\emph{Parties \uncertain{saturées}} de $|F|$ : p. ex. les $X \in F$.

\struck{Définition — Soient $F$, $G$ deux figures. On dit qu'elles sont \emph{compatibles}, si les conditions équivalentes : $|F| \cap |G|$ est saturé pour les relations d'équivalence sur $|F|$, $|G|$, et la rel.}
\note{ces lignes sont barrées d'un trait horizontal par ligne et de deux traits obliques}

\page{27}
\emph{Cont 3} \quad Si $F$ est une \struck{\ill{}} figure, et $F'$ une partie fermée de l'ens. ordonné $F$ (i.e. $X, Y \in F$, $Y \in F'$, $X \leq Y \Rightarrow X \in F'$) alors $F'$ est une figure.

On dit que $F'$ est une \emph{sous-figure} de $F$ — ainsi, les sous-figures de $F$ sont en corr. biunivoque avec les parties fermées de l'ens. ordonné $F$.

\uncertain{Attention}, si $F'$, $F$ sont deux figures, $F' \subset F$ \ill{} que $F'$ soit une sous-figure. (\ill{} l'exemple-type, prenant $F$ formé des $X_0 \subset X_1$, $X_1$ une variété connexe, $F' = \{X_1\}$.)

On écrira \struck{$X$} $F' \leq F$ pour $F'$ sous-figure de $F$. C'est une relation d'ordre entre figures.

\struck{Soient $F$, $G$ deux figures, \ill{}. On dit que $F$ et $G$ sont \emph{compatibles}, si $F \cup G$ est}
\note{ces lignes sont hachurées ; une note oblique dans la marge gauche, en regard, n'est pas lue}

\page{28}
Soient $F$ et $G$ deux figures. Conditions équivalentes (où $F$ et $G$ sont des ens. \emph{admissibles} de parties de $L$) :
\begin{enumerate}
  \item[a)] $|F| \cap |G|$ est une partie saturée de $|F|$, et de $|G|$, et les deux rel. d'équiv. induites sur elle sont égales.
  \item[b)] $\forall$ strates ouvertes $U$ ($= X^{\circ}$) de $F$ et $V$ de $G$, on a $U \cap V = \emptyset$ ou $U = V$.
  \item[b')] \struck{\ill{}} $X^{\circ} \cap Y^{\circ} = \emptyset$ ou $X = Y$.
  \item[c)] $\forall X \in F$, $Y \in G$, $X \cap Y$ est saturé dans $|F|$ et dans $|G|$ (pour la saturation \ill{} par rapport aux str. \ill{} des multiplicités $\widehat{X}$, $\widehat{Y}$).
  \item[d)] $F \cup G$ est \uncertain{admissible}, et $F$ et $G$ sont des \ill{} \emph{fermés}.
\end{enumerate}
\note{la ligne b') n'est faite que de tirets suivis de la formule}

On dit alors que $F$ et $G$ sont « \ill{} compatibles ». Il en est ainsi si $F$ et $G$ sont des sous-figures d'une même figure $H$ \ill{} : $F \cap G$ fermé dans $F$ et dans $G$, $F$, $G$ fermés dans $F \cup G$, \ill{} ordonné $\simeq F \amalg_{F \cap G} G$.
\note{cette ligne se mêle à la note de marge ; lecture fragmentaire}

Si $F$ et $G$ sont \emph{compatibles}, [s']ils le sont ensemblistement, et $F \cup G$ est une \ill{} ; $F$ et $G$ sont des sous-figures de $F \cup C$.
\note{ainsi sur la page : « $F \cup C$ » pour $F \cup G$}

Si $F$, $G$ sont des sous-figures d'une figure $H$, alors ils \ill{} ensemblistement, et compatibles (car $F \cup G$ est une \ill{} de $H$, donc une figure).

\uncertain{Configuration} : ensemble de figures deux à deux compatibles.

\marginal{N.B. $F$ et $G$ \ill{}}
\note{deux longues notes obliques dans la marge gauche, ouvertes par « NB », mentionnent $F \cap G$, $F \cup G$, « Ex. » et « compatibles » ; le reste n'est pas lu}

\page{29}
\emph{Cont 4} \quad Si $F$, $G$ \struck{$H$ sont des figures} sont \add{deux} figures, telles que toute multistrate de $F$ est comp. à toute multistrate de $G$*, alors $F$ et $G$ sont compatibles (l'inverse étant évident).
\marginal{*) $F$ et $G$ ens. comp.}

\struck{Cor. Si $F$, $G$, $H$ sont des figures \ill{}} \struck{Équiv.}

\emph{Cor.} Soit $(F_i)_{i \in I}$ une famille \struck{\ill{}} de figures. Pour que les $F_i$ soient \ill{} en configuration, il faut et il suffit que pour $i \neq j$ et toute multistrate $\widehat{X}$ de $F_i$, $\widehat{Y}$ de $F_j$, $\widehat{X}$ et $\widehat{Y}$ soient compatibles. Si $I$ est fini, cela implique que $F = \bigcup F_i$ est une figure, et $F_i$ en est une des sous-figures.
\note{l'énoncé du corollaire est très surchargé ; la lecture donnée suit les formules, sûres, et reste douteuse pour les mots qui les relient}

\struck{N.B. Inversement, \ill{} dans le cas d'une paire $\{F, G\}$ de figures (à savoir, \ill{})}
\note{ces lignes sont barrées de traits obliques}

\emph{Contrées \struck{ensemblistes} \add{de points}} \quad La contrée est dite « ensembliste », ou « contrée de points » (\uncertain{c'est-à-dire} les lieux prennent le nom de \emph{points}) si deux figures ensemblistement compatibles sont compatibles (ce qui rend donc tautologique \emph{Cont 4}).

\page{30}
Pour les contrées non ensemblistes, il y a lieu de supposer ceci (à vérifier sur les exemples \uncertain{visuels} non ensemblistes, à suivre)

\emph{Cont$'_4$} \quad Si $F$, $G$ sont deux figures \struck{\ill{}} telles que \ill{} \ill{} compatibles, \ill{} qu'elles le soient ensemblistement, et que pour deux strates ouvertes $U$ de $F$, $V$ de $G$, \struck{on ait} $U \neq V$, et pour \struck{tout} $x \in U$, $y \in V$, \struck{$\{x, y\}$} $\{\{x\}, \{y\}\}$ soit une figure (i.e. $\{\{x\}\}$ et $\{\{y\}\}$ soient compatibles).
\note{le numéro est surchargé : un « 4 » et un « 1 » se lisent l'un sur l'autre après « Cont », avec un accent ; l'énoncé est très surchargé, une insertion interlinéaire oblique (« toute multistrate $\widehat{X}$ de $F$ est comp. à toute $\widehat{Y}$… ») n'est lue qu'en partie, et la note de marge qui la prolonge (« si $x \in U$, $y \in V$… disjoints ») aussi}

N.B. Pour tout $x \in L$, \ill{} \struck{\ill{}} \add{\ill{}} la seule strate \struck{\ill{}}. Si $x, y \in L$, disons que $F_x$, $F_y$ sont compatibles signifie que $\{F_x, F_y\}$ est \struck{\ill{}} une figure (ce qui \uncertain{toujours} vrai si $x = y$). Si $x \neq y$, \struck{dites} on dit que $x$, $y$ sont « \emph{disjoints} » s'ils sont distincts et compatibles, i.e. \ill{} alors que $\{x, y\} \in \mathfrak{P}_2(L)$ est une $0$-sphère de $C$.

Deux parties $A$, $B$ de $L$ sont dites \emph{disjointes} (\ill{} contrée $C$) — ou \emph{$C$-disjointes} — si $\forall x \in A$, $\forall y \in B$, $x$, $y$ sont disjoints ($\Rightarrow A \cap B = \emptyset$).

\page{31}
Ainsi \emph{Cont$'_4$} dit que deux figures $F$, $G$ sont compatibles ssi pour $U$ strate ouverte de $F$ et $V$ de $G$, \struck{telles que $U \neq V$, on ait} on ait $U = V$ ou $U$ et $V$ sont $C$-disjoints.

\subsection{12 juin}
\note{date de sa main, « 12. Juin », soulignée, en tête de l'alinéa}

Depuis hier j'ai pas mal réfléchi à la construction des pseudo-disques topologiques. Il est apparu que si \struck{\ill{}} $(D, \partial D)$ est un tel pseudo-disque (ou \ill{}), les « lieux » sur $D - \partial D$ ne \uncertain{sont} \struck{\ill{}} \ill{} seulement des parties (compactes contractiles etc.) de $D - \partial D$, satisfaisant à certaines conditions, mais elles doivent être accompagnées d'une donnée supplémentaire, \ill{} \ill{} \ill{} la frontière \ill{} $\partial D$. Du coup, \ill{} que l'interprétation \struck{\ill{}} que j'avais \ill{}, d'une structure de contrée sur \ill{} donnée de l'ens. $L$ des lieux, et de certains ensembles de parties de $L$, les « figures », ne \ill{} pas
\note{la seconde moitié de l'alinéa est écrite vite ; seuls les termes mathématiques sont sûrs}

\page{32}
valable pour $C$ « topologique combinatoire ». Donc il faut prendre plutôt :
\[
  C \ \text{donnée de}\ C = (\mathcal{M}, \ll, R_{\mathcal{M}}, \mathfrak{F} \ \text{ens. de parties de}\ \mathcal{M})
\]
\note{sous $\ll$ est écrit « relation d'ordre » ou « raffinement » (lecture douteuse), sous $R_{\mathcal{M}}$ « la compatibilité, rel. sym. réfl. » ; la suite de la ligne est barrée}
\marginal{$R_{\mathcal{M}}$ résultant de $\ll$, $\mathfrak{F}$ : $X, Y \in \mathcal{M}$ sont \ill{} $X, Y \in F$ \ill{}}
\note{note oblique de la marge gauche, lue en partie}

Plutôt que de dégager \ill{} l'axiomatique, dégager une terminologie géométrique \ill{} en contrées \ill{} ; les axiomes que je \ill{} au passage ; l'axiomatisation \ill{} ensuite.
\note{deux lignes très rapides ; la lecture n'en retient que l'intention}

\emph{Sous-contrée} : est définie par un sous-ensemble $\mathcal{M}'$ \struck{\ill{}} (qui est fermé \ill{}) dans $\mathcal{M}$ pour $\ll$, i.e. si $X$ est dans $\mathcal{M}'$, tout raffinement de $X$ y est aussi. \struck{Et} On munit $\mathcal{M}'$ des relations induites $\ll$ et $R_{\mathcal{M}'}$ \struck{\ill{}} par $\ll_{\mathcal{M}}$ et $R_{\mathcal{M}}$. Mais \ill{} qu'il n'est pas \struck{évident} \add{dit} que les axiomes qu'on va imposer à une contrée soient stables par passage à une
\note{une note oblique au bas de la marge gauche, ouverte par « NB », n'est pas lue}

\page{33}
sous-contrée ; \ill{} \ill{}. \ill{} nous réservons le nom de « contrée » à \ill{} systèmes $(\mathcal{M}, \ll, R_{\mathcal{M}})$ \ill{}. Je vais \ill{} \ill{} axiomes, particulièrement cruciaux.

\emph{Cont 1} \quad La relation « $X \leq Y$ » ($\Leftrightarrow X \ll Y$ et $X$ comp. $Y$) est une relation d'ordre.
\note{double trait vertical en marge}

\struck{Cont 2} (\emph{Fig 2}) \struck{\ill{}} Si $F \in \mathfrak{F}$, deux éléments $X$, $Y$ de $F$ sont compatibles, et $F$ est une partie fermée de $\mathcal{M}$ ($\leq$).

\struck{Cont} (\emph{Fig 3}) Si $F$ \struck{et $G$ sont des} \add{est une} figure, \ill{} partie $F'$ de $F$ qui est fermée, est une sous-figure \ill{}. (La partie vide de $\mathcal{M}$ est une figure (figure vide).)
\note{la parenthèse finale est encadrée ; le numéro barré devant « Fig 3 » est illisible}

\struck{Cont} (\emph{Fig 1}) Si $X \in \mathcal{M}$, l'ens. $F_X = \mathcal{M}_{\leq X}$ des $Y \in X$ \struck{tels} est une figure, qui s'appelle la « \emph{figure élémentaire} » définie par $X$ — la multiplicité $X$.
\note{ainsi sur la page : « $Y \in X$ » pour $Y \leq X$}
\marginal{(de Cont 2 et Cont 4)}

\emph{Cor.} (si $\mathcal{M} \neq \emptyset$) \struck{\ill{}} Fig 3 donne \ill{} vide —

N.B. Si $\mathcal{M} = \emptyset$ \ill{} la seule partie de $\mathcal{M}$ (\uncertain{à savoir} la partie vide) est une figure.

\struck{Cor. (de Fig 1 et Fig 2) Si $Y, Z \ll X$, et si $Y$, $Z$ sont comp. à $X$, \ill{} sont compatibles}
\note{cette dernière ligne est barrée de traits obliques. Deux longues notes obliques de la marge gauche (sur les sous-figures $F' \subset F$ et sur $X \leq Y$, $Y \leq X$, « $X$ comp. $Y$ ») ne sont lues qu'en fragments}

\page{34}
\struck{Cont \ill{}}
(\emph{Fig 4}) Soient $F$ et $G$ deux figures, telles que pour \struck{$X$} $X \in F$, $Y \in G$, $X$ et $Y$ comp. Alors $F \cup G$ est une figure — (\emph{figure réunion}).
\note{« Fig 4 » est entouré, écrit sur un numéro « Cont » barré ; « figure réunion » est doublement souligné}

Bien sûr $F \cap G$ est une sous-figure de $F$ et de $G$.

\struck{Corollaire (de Fig 1, Fig 2, Fig 4 \ill{}) Si $X$, $Y$, $X'$, $Y'$ sont tels que $X' \ll X$, $Y' \ll Y$, et \ill{} comp., alors $X'$, $Y'$ comp.}

\emph{Contrées finitistes} \struck{Multiplicités finies :}

\emph{Corollaire} (de Cont 2, Cont 3, Cont 5) Les figures \emph{finies} \struck{\ill{}} \ill{} les parties finies de $\mathcal{M}$, \emph{fermées} pour $\leq$, formées de multiplicités deux à deux \emph{compatibles}.

\struck{Donc} \struck{les structures} \struck{les figures finies \ill{}} Donc une contrée finitiste est connue, quand on connaît seulement ses structures $\ll$, $R_{\mathcal{M}}$ (ordre, et rel. sym. réfl.). \struck{et ce}

À quelles conditions \uncertain{sur} deux telles relations, satisfaisant à Cont 1, les parties finies $S$ fermées \struck{\ill{}} et formées d'éléments mut. compatibles (i.e. les \emph{configurations}) satisfont-elles Cont 2 à Cont 5 ?
\note{une note oblique de la marge gauche (« … dès que $F$ et $G$ sont deux figures compat. … ») n'est lue qu'en partie}

\page{35}
\struck{\ill{}} La condition \ill{} est la suivante ; \struck{\ill{}}

\emph{Cont$'_3$} (F. multiplicités dans le cas non finitiste) \struck{\ill{}} Pour tout $X \in \mathcal{M}$, \struck{la figure} l'ens. $\mathcal{M}_{\leq X}$ est fini. (i.e. les figures élémentaires sont finies)
\note{le « 3 » en indice est écrit sous un « 4 » barré}

Pour contrées, on exige un peu plus : que la \emph{dimension combinatoire} de tout $X \in \mathcal{M}$ soit finie (i.e. $\dim \mathrm{comb}\, \mathcal{M}_{\leq X}$ finie), et plus précisément encore

\emph{Cont 3} \quad Si $X \in \mathcal{M}$,
\[
  \underbrace{\sup_{X' \ll X} \dim \mathrm{comb}\, \mathcal{M}_{\leq X'}}_{\text{dimension (géom.) de } X} < \infty
\]
\note{triple trait vertical dans la marge, et, à côté, un mot oblique doublement souligné (« demain » ?) non assuré}

On appelle \emph{lieu} un élément de $\mathcal{M}$ minimal pour $\ll$. On désigne par $\mathcal{L} \subset \mathcal{M}$ l'ens. des lieux.

Pour tout $X$, son \emph{ombre} (\add{fermée}) sur $\mathcal{L}$ est l'ens. $\mathcal{L}_X$ des lieux $x$ \struck{\ill{}} $x \ll X$. On \ill{} « \ill{} de $X$ » i.e. \ill{} \ill{}. On pose
\[
  \mathcal{L}^{\circ}_X = \mathcal{L}_X \smallsetminus \bigcup_{Y < X} \mathcal{L}_Y .
\]
\marginal{$\mathcal{L}_F = \bigcup_{X \in F} \mathcal{L}_X$ \quad ombre de $F$}
\note{la note de marge est oblique ; une seconde note oblique, au-dessus (« … fermée », « … lieux »), n'est pas lue}

\emph{Cont 4} \quad a) $\forall X \in \mathcal{M}$, on a $\mathcal{L}^{\circ}_X \neq \emptyset$

b) Si $X$, $Y$ sont compatibles et distincts, $\mathcal{L}^{\circ}_X \cap \mathcal{L}^{\circ}_Y = \emptyset$.
\note{le « 4 » est écrit sur un autre chiffre}

\page{36}
\emph{Scholie} \struck{L'applicat.} Soit $F$ une figure. Alors

a) L'application
\[
  X \longmapsto \mathcal{L}_X = \widetilde{X}, \qquad F \longrightarrow \mathfrak{P}(\mathcal{L})
\]
est injective, et son image $\widetilde{F}$ est un ensemble de parties « admissible » de l'ensemble $\mathcal{L}$, i.e.
\begin{enumerate}
  \item[a)] $\forall \widetilde{X}, \widetilde{Y} \in \widetilde{F}$, on a $\widetilde{X} \cap \widetilde{Y} = \bigcup_{\widetilde{Z} \in \widetilde{F},\ \widetilde{Z} \subset \widetilde{X} \cap \widetilde{Y}} \widetilde{Z}$ ;
  \item[b)] $\forall \widetilde{X} \in \widetilde{F}$, posant $\partial \widetilde{X} = \bigcup_{\widetilde{Y} \in \widetilde{F},\ \widetilde{Y} \subsetneq \widetilde{X}} \widetilde{Y}$, et $\widetilde{X}^{\circ} = \widetilde{X} \smallsetminus \partial \widetilde{X}$, on a $\widetilde{X}^{\circ} \neq \emptyset$.
\end{enumerate}
On \struck{b) L'appl. \ill{}} \ill{} $\widetilde{X}^{\circ}$ l'\emph{ombre ouverte} de $X$ sur $\mathcal{L}$, $\partial \widetilde{X}$ l'\struck{ombre fr.} \add{\emph{bord d'ombre}} de $X$.

b) L'application
\[
  X \longmapsto \mathcal{L}^{\circ}_X = \widetilde{X}^{\circ}, \qquad F \longrightarrow \mathfrak{P}(\mathcal{L})
\]
est injective, et \struck{\ill{}} les images forment une \emph{partition} de l'ombre $\mathcal{L}_F$ de la figure $F$, \uncertain{indexée} par $F$.

c) \struck{Soit} La relation $\leq$ sur $F$ est celle induite par la relation d'inclusion sur les ombres. Si $X, Y \in F$, les conditions équiv. :
(i) $X \leq Y$ \quad (ii) $\widetilde{X} \subset \widetilde{Y}$ \quad (iii) $\widetilde{X}^{\circ} \subset \widetilde{Y}$ \quad (iv) $\widetilde{X}^{\circ} \cap \widetilde{Y} \neq \emptyset$.
\note{les conditions (i)–(iv) sont réunies par une accolade}
\note{« bord d'ombre » est doublement souligné, écrit au-dessus de mots barrés (« ombre fr. », lecture douteuse)}
\marginal{Déf. \ill{} $\partial X = \lbrace Y \in \mathcal{M} \mid Y < X \rbrace$ \ill{} celle \ill{} le bord de $X$}
\note{note oblique de la marge gauche, en regard de b) ; le premier mot, « Déf. », est barré en surcharge}

\page{37}
\begin{enumerate}
  \item[d)] \struck{Les} Pour toute sous-figure $F'$ de $F$, on a $\widetilde{F}' = \mathcal{L}_{F'} \subset \mathcal{L}$ son ombre, et
  \[
    F' \longmapsto \widetilde{F}'
  \]
  est une bijection de l'ens. des sous-figures de $F$, sur l'ens. des parties (de $\widetilde{F}$ (ombre de $F$)) qui sont « fermées » \ill{} pour la rel. d'admissibilité de $\widetilde{F}$ par les $\widetilde{X}$ ($X \in F$).
\end{enumerate}
\note{l'écriture de la première formule est surchargée ($\mathcal{L}_F$ récrit en $\widetilde{F}'$) ; lecture douteuse}

\emph{Proposition.} Soit $\mathcal{M}$ une contrée, $\mathcal{M}'$ une sous-contrée, \struck{\ill{}} i.e. une partie de $\mathcal{M}$, fermée pour $\ll$, et munissons $\mathcal{M}'$ des relations $\ll_{\mathcal{M}'}$, $R_{\mathcal{M}'}$ induites \ill{}, et de l'ens. $F_{\mathcal{M}'}$ des parties \struck{\ill{}} de $\mathcal{M}'$ qui sont des figures de $\mathcal{M}$. Alors
\[
  (\mathcal{M}', \ll_{\mathcal{M}'}, F_{\mathcal{M}'})
\]
satisfait les conditions \emph{Cont 1} à \emph{Cont 4}, et si $X \in \mathcal{M}'$, ses dimensions combinatoire et géométrique sont les mêmes dans $\mathcal{M}'$ et dans $\mathcal{M}$ ; si $\mathcal{L}' = \mathcal{M}' \cap \mathcal{L}$, alors pour $F \in F_{\mathcal{M}'}$, son ombre (\add{fermée}) dans $\mathcal{M}'$ et dans $\mathcal{M}$ sont les mêmes ; si $X \in \mathcal{M}'$, \struck{\ill{}} son ombre ouverte est la même \add{pour $\mathcal{M}'$ et pour $\mathcal{M}$}, \ill{} le bord $\partial X$. Les conditions \emph{Fig 1} à \emph{Fig 4} sont satisfaites dans $\mathcal{M}'$. La condition \emph{Cont$'_3$} \struck{\ill{}}

\page{38}
également satisfaite par passage à une sous-contrée. Enfin, si $C$ est finitiste, $C'$ aussi.

\emph{\struck{Figures} disjointes} \quad On dit que deux \ill{} sont disjointes si elles sont compatibles, et si leur intersection \ill{} vide.
\note{au-dessus de « Figures », barré, une insertion oblique en petits caractères n'est pas lue}
\ill{} \emph{multistrates} \ill{}. On dit que $X$, $Y$ sont \emph{disjointes} si \struck{\ill{}} \ill{} $F_Y$ \ill{}, i.e. s'il n'existe pas de \struck{\ill{}} $Z \in \mathcal{M}$, avec $Z \leq X, Y$.

\struck{\ill{}} donné deux figures $F$, $G$, on dit que $F$ \emph{raffine} $G$ (et on écrit $F \lll G$) si $\forall X \in F$, $\exists Y \in G$ tel que $X \lll Y$. Pour deux figures élémentaires $F_X$, $F_Y$, cela signifie donc simplement $X \ll Y$.
\note{« raffine » est triplement souligné ; le signe entre figures est un triple chevron, entre $X$ et $Y$ ensuite aussi}

\struck{Si $F$ raffine $G$}

\emph{Proposition.} Pour que $F$ soit une sous-figure de $G$, il faut et il suffit que $F$ raffine $G$, et soit compatible à $G$.
\note{un trait vertical en marge de l'énoncé}

\ill{} $\Rightarrow$ \uncertain{trivial}, pour $\Leftarrow$ on note que $X \ll Y$ \struck{et} $X$ comp. à $Y$ \struck{\ill{}} implique $X \leq Y$, donc $X \in G$ à cause de \emph{Fig 2}.
\note{la marge gauche porte une longue note oblique, reliée au texte par un trait, sur les figures disjointes ($\forall X \in F$, $Y \in G$, « disj. ») ; elle n'est lue qu'en fragments}

\page{39}
\emph{Cont 5} \quad Soient $X, Y \in \mathcal{M}$, $X$ et $Y$ \emph{disjoints}. Alors pour $X' \ll X$ et $Y' \ll Y$, $X'$ et $Y'$ sont \struck{disjoints} \struck{compatibles} disjoints.
\marginal{\ill{} si \struck{Cont 2} \ill{} (voir 33)}
\note{double trait vertical en marge ; « voir 33 » renvoie à la page 33}

\struck{Cela signifie donc 1°) que $X'$ et $Y'$ sont compatibles et 2°) qu'il n'existe pas $Z \leq X'$, $Z \leq Y'$. Alors $X'$, $Y'$ sont disjoints, car si $Z \in \mathcal{M}$ avec $Z \leq X'$, $Z \leq Y'$, alors $Z \leq X$, \ill{}}
\note{ce passage est encadré puis barré de traits obliques}

Sa signification est donc : \struck{\ill{}} (a) toutes les conditions des \add{deux} \emph{compatibles}, $X'$ et $Y'$ sont \emph{compatibles} ; (b) \struck{\ill{}} $\nexists\, Z \in \mathcal{M}$ avec $Z \ll X$, $Z \ll Y$.
\note{(a) et (b) sont entourés}

(i.e. si $X$, $Y$ compatibles, dans $\mathcal{M}_{\leq X} \cap \mathcal{M}_{\leq Y} = \emptyset$ implique $\mathcal{M}_{\ll X} \cap \mathcal{M}_{\ll Y} = \emptyset$)

\emph{Cor.} Si $F$ et $G$ sont des figures disjointes, deux raffinements $F'$ de $F$, $G'$ de $G$, $F'$ et $G'$ sont disjointes.

\emph{Remarque.} Soit $\mathcal{M}'$ une sous-contrée \ill{} ; deux multistrates $X$, $Y$ \struck{\ill{}} \ill{} deux figures $F$, $G$ de $\mathcal{M}'$, sont disjointes \ill{} de $\mathcal{M}'$ ssi elles le sont dans $\mathcal{M}$ \ill{}, et donc la condition \ill{} (itou pour \emph{Cont 6} qui suit) : si $\mathcal{M}$ satisfait \emph{Cont 5}, $\mathcal{M}'$ aussi.

\emph{Cont 6} (F. \uncertain{cellulaire} ?) \quad Si $X, Y \in \mathcal{M}$ sont tels que \struck{$\forall$} $\widetilde{X} \cap \widetilde{Y} = \emptyset$ \emph{et} tout \struck{$x \in$} $x \in \widetilde{X}$ est comp. à tout $y \in \widetilde{Y}$, alors $X$ et $Y$ sont disjoints.
\marginal{donc $\forall x \in \widetilde{X}$ disjoint de tout $y \in \widetilde{Y}$}
\note{triple trait vertical en marge de Cont 6 ; deux autres notes obliques de la marge gauche, reliées au texte par des flèches, ne sont lues qu'en fragments}

\page{40}
\emph{Contrées induites}

1°) Soit $F$ une figure. Contrée $C_F$ \add{induite}, \struck{définie par les} \add{i.e. la} partie fermée de $\mathcal{M}_{\ll}$ engendrée par $F$, i.e. l'ens. des $X \in \mathcal{M}$ tels que $X$ raffine $F$.

\emph{Contrée \uncertain{résiduelle}} \quad $C_{\complement F}$ la \add{sous-}contrée formée des $X$ \add{$\in \mathcal{M}$} qui sont disjoints de $F$. Par \emph{Cont 5}, \ill{} est \struck{\ill{}} \add{fermée} dans $\mathcal{M}^{\sharp}_{\ll}$.
\note{l'exposant de $\mathcal{M}$, lu $\sharp$, est douteux}

\struck{Cor.} Soient $F$ une figure, $G$ une \ill{}. \struck{Subdivision} La contrée résiduelle de $F$ dans $G$ est par définition
\[
  C_F \cap C_{\complement G}
\]
\note{les lettres $F$ et $G$ de la formule sont surchargées}

\emph{Subdivision} \quad Soit $F$ \add{$\subset \mathcal{M}$} une figure. On appelle \emph{subdivision} de $F$ une figure $F' \subset \mathcal{M}$ telle que
\begin{enumerate}
  \item[a)] $F'$ soit un raffinement de $F$
  \item[b)] $F'$ soit maximale, \add{pour $\subset$}, parmi les \struck{raffi} figures qui raffinent $F$.
\end{enumerate}

N.B. Si $F'$ raffine $F$, alors son ombre $\widetilde{F}'$ est contenue dans celle $\widetilde{F}$ de $F$.

\marginal{N.B. L'ensemble des lieux de la contrée $C_F$ \ill{} l'ombre \ill{} de $F$}
\note{une seconde note oblique de la marge gauche, sur l'ensemble des lieux de $C_{\complement F}$, n'est lue qu'en fragments}

\end{document}
